\subsection{Sources, comparison maps, and the two remaining image problems}

We use the generators and positive tangential normalization of the task.
The source question is LM7 D.13, developing the comparison of groups in
\emph{La Longue Marche}, \S51.15 (cote 140-4, scans 156--170).
Only the supplied transcription, not the underlying scans, was available
for this comparison. We make no assertion of a new reading of those scans.

Put $F=\langle x,y\rangle$, let $\widehat F$ be its profinite completion,
and write $F_\ell$ for its maximal pro-$\ell$ quotient. The quotient
$q_\ell:\widehat F\twoheadrightarrow F_\ell$ has characteristic kernel.
It therefore induces a continuous homomorphism on outer automorphism
groups. Let $I^n$ denote the filtration obtained by acting trivially,
up to inner automorphism, on $F_\ell/\Gamma_{n+1}F_\ell$.
We filter each full profinite source by inverse images of $I^n$.
The difference of an automorphism and the identity, read in
$\Gamma_{n+1}/\Gamma_{n+2}$, is a derivation. Composition becomes addition
in this quotient, and conjugating the automorphism changes that
derivation by an inner derivation. Thus we obtain the specified maps
\[
 (I^nG/I^{n+1}G)\otimes_{\Z_\ell}\Q_\ell
 \longrightarrow \operatorname{OutDer}_{\mathrm{sp}}(L_\ell)_n
 \longrightarrow \operatorname{OutDer}(L_\ell/L_\ell'')_n,
 \qquad L_\ell=\operatorname{Lie}_{\Q_\ell}(x,y).
\]
Here $G$ is one of the full profinite source groups, and the first
tensor product means the rational span of its image in the integral
graded automorphism module. Expanding each peripheral conjugacy equation
gives the special-derivation equation. Expanding the commuting equations
with the geometric $S_3$ gives invariance. These are necessary conditions
on the images of the displayed maps.

\begin{statement}{D.1}{Open}{The weight-five image of the full profinite special group}
For $\ell\geq5$, the dimension requested in Task D.1 has not been
determined here. The target computed in R8 is
\[
 (X+Y)\,\langle X^2,XY,Y^2\rangle_{\Q_\ell}[x,y].
\]
Statement D.21 proves that this image has dimension either $1$ or $3$:
the actual geometric $S_3$ preserves it, and its invariant gamma line
is already present by D.20. The remaining choice between $1$ and $3$
is not determined here.
\end{statement}

\begin{statement}{D.2}{Open}{The extra weight-nine direction for the full profinite gentle group}
For $\ell\geq5$, we have neither constructed an element of
$\widehat{\mathrm{GT}}_0$ with leading metabelian coordinate outside
$\Q_\ell G_9$, nor proved that every such coordinate belongs to that line,
where
\[
 G_n=\frac{(X+Y)^n-X^n-Y^n}{XY},\qquad
 p_{\mathrm{extra}}=X^2Y^2(X+Y)^3.
\]
The requested choice between dimensions $1$ and $2$ remains open in
this return.
\end{statement}

\begin{statement}{D.3}{Proved}{The finite pentagon comparison in the weight-nine stratum}
At the compatible quotients
$PB_n/(\gamma_{11}PB_n\,PB_n^{101})$, for $n=3,4$, the complete
multiplier-one, weight-at-least-nine gentle shadow stratum has
$101^{37}$ elements. Exactly $101^2$ of these satisfy the complete
pentagon. Their degree-nine leading spaces have dimensions four and
one, respectively; their metabelian images have dimensions two and
one. The latter image is the line generated by $G_9$.
Statements D.7--D.8 give the finite comparison and an exact integer
certificate. Its characteristic-zero comparison, together with the
actual Galois lower bound in D.20, derives R9 in weight nine.
The enumeration is of the entire stratum that can contribute a
weight-nine leading term, with its degree-ten corrections; it is
not an enumeration of other multipliers or of lower leading weights.
\end{statement}

\subsection{What the linear and literature checks establish}

\begin{statement}{D.4}{Proved}{Recomputed targets and the checked pentagon upper bound}
The original rational program reproduces R8's four rows in weights
$1,\ldots,10$. In weight nine its invariant metabelian image has basis
$p_0=p_{\mathrm{extra}}$ and
\[
\begin{split}
 p_1={}&2X^7+8X^6Y+17X^5Y^2+23X^4Y^3\\
       &+23X^3Y^4+17X^2Y^5+8XY^6+2Y^7,
\end{split}
\qquad G_9=\frac{15}{2}p_0+\frac92p_1.
\]
The leading metabelian image of full $\widehat{\mathrm{GT}}$ is contained
in $\Q_\ell G_n$ in positive odd weight $n$, and is zero in positive
even weight. In particular $p_{\mathrm{extra}}$ is excluded in weight
nine by this implication from the pentagon.
\end{statement}
\begin{proof}
We reran the supplied programs, retaining their provenance under
\texttt{code/de\_source/}; this is a reproducibility check rather than
an independently implemented rational Lie calculation. Their complete
exact rational outputs and source witnesses are included. The rows are
\[
\begin{array}{c|rrrrrrrrrr}
 n&1&2&3&4&5&6&7&8&9&10\\\hline
 \mathrm{special}&0&0&1&0&3&0&6&4&13&12\\
 \mathrm{metabelian\ image}&0&0&1&0&3&0&5&0&7&0\\
 S_3\text{-}\mathrm{invariants}&0&0&1&0&1&0&2&1&4&2\\
 \mathrm{their\ image}&0&0&1&0&1&0&1&0&2&0.
\end{array}
\]
The calculation solves $[f,y]=[g,x+y]$ in the free Lie algebra,
then applies the canonical quotient $L\to L/L''$ to its solutions.
This order of operations specifies the image being counted.

For the upper bound we use the opened primary text of Furusho,
Corollary 0.3, Theorem 0.5, Proposition 5.1 and Lemma 5.2
\cite{DE_Furusho2011}. The leading homogeneous term of a profinite
pentagon equation, after applying $q_\ell$ and the continuous rational
unipotent completion, satisfies the infinitesimal pentagon: multiplication
of factors of the same least positive degree becomes addition, and
the standard braid insertion maps induce the corresponding insertion
maps on the graded Lie algebras. Thus this leading term lies in
$\mathfrak{grt}_1(\Q_\ell)$. Furusho's result implies its gamma identity.
We spell out the coordinate conversion. In the completed tensor algebra
write $h=h_xx+h_yy$, and put $M(h)=(h_yy)^{\mathrm{ab}}$.
Expansion gives
\[
 M([x,y])=XY,\quad M([x,h])=XM(h),\quad
 M([y,h])=YM(h)\quad(h\in L').
\]
Moreover $M([L',L'])=0$, because the commutative images of both
derived factors vanish. Hence $M(p[x,y])=XYp$ on $L'/L''$.
The homogeneous gamma identity reads
$M(h)=a_n(X^n+Y^n-(X+Y)^n)$. Cancellation of $XY$ gives the asserted
line. The special equation implies $(X+Y)\mid p$: in $L'/L''$ it
reads $-Yp=-(X+Y)q$, and $X+Y$ is relatively prime to $Y$.
For even $n$, evaluating the numerator at $Y=-X$ makes $a_n=0$.
Weight one is zero directly. The rerun gamma verifier checks this Fox
coordinate identity on all 224 Lyndon words in degrees two through ten
and the displayed independent weight-nine vectors.
\end{proof}

The lower bound in R9 is established in D.20 from the opened primary
Ihara formula, with its exact Fox coordinate and an extraction inside
the actual compact Galois image. The initial Hain--Matsumoto audit
reached Theorem 8.4, \S8.4, Lemma 8.5 and Remark 8.6
\cite{DE_HainMatsumoto2003}; that remark explicitly points to Ihara
for nontriviality. The subsequent primary-source check opens Ihara's
ICM lecture, \S6.2, the theorem of \S6.3, and (6.4.4), and supplies
the missing direct comparison in D.20. Brown's motivic formula
\cite[(1.3)]{DE_Brown2015} is not needed for that arithmetic proof.
The source program's field
\texttt{actual\_Galois\_and\_GT\_rational\_rank} remains a
literature-populated label; the arithmetic conclusion follows from
the proof in D.20 rather than from that label.

\begin{statement}{D.5}{Known}{Cyclic-cover compatibility for every full profinite GT element}
Let $H_N=\ker(\widehat F\to\Z/N)$, where $x\mapsto1$ and $y\mapsto0$.
The continuous homomorphism
\[
 \delta_N:H_N\longrightarrow\widehat F,
 \quad x^N\mapsto x,\quad y\mapsto y,\quad
 x^ayx^{-a}\mapsto1\quad(0<a<N)
\]
is defined on a profinite free basis. For every
$(\lambda,f)\in\widehat{\mathrm{GT}}$, with the convention
$x\mapsto x^\lambda$, $y\mapsto f^{-1}y^\lambda f$, one has
\[
 \delta_N(f)=y^{c_N(f)}f,\qquad
 c_{MN}(f)=c_M(f)+c_N(f),
\]
where $c_N(f)$ is the $y$ coordinate of $\delta_N(f)$ in
$\widehat F^{\mathrm{ab}}$. These are Enriquez's Theorem 8.1 and
Remark 8.6 \cite{DE_Enriquez2007}.
\end{statement}
\begin{proof}[Comparison with the covering map in LM7]
For the negative Schreier generator $e_j=x^{-j}yx^j$, $0<j<N$,
\[
 e_j=x^{-N}(x^{N-j}yx^{-(N-j)})x^N.
\]
Both sides are in $H_N$, and applying $\delta_N$ gives $1$.
The maps denoted $\phi_N$ in LM7 D.20 and $\delta_N$ above consequently
agree on every member of a free basis, hence agree as continuous
homomorphisms. We retain $\rho_N^{\rm dist}:=-c_N$ as
the explicitly defined negative distribution coordinate. The earlier
sentence identifying it with the positive-root Kummer character in
the displayed Galois convention requires the following correction.
For $N=2$, the exact comparison is
\[
 c_2(f_\sigma)=\kappa_2(\sigma),\qquad
 \rho_2^{\rm dist}(f_\sigma)=-\kappa_2(\sigma).
 \tag{D.5.1}
\]
Here $\kappa_2$ is defined by
$\sigma(2^{1/m})=\zeta_m^{\kappa_2(\sigma)}2^{1/m}$
using positive roots and $\zeta_m=e^{2\pi i/m}$.
We now prove the coefficient and convention maps.

The continuous homomorphism $c_2:H_2\to\widehat\Z$ is the
$y$ coordinate after $\delta_2$. On $\widehat F'$ it kills
$\widehat F''$. In the free basis
$x^2,y,xyx^{-1}$ of $H_2$, every element of $\widehat F'$
has zero first coordinate and its other two abelian coordinates
sum to zero. Conjugation by $x$ interchanges those last two
coordinates, and conjugation by $y\in H_2$ fixes them.
Moreover $[x,y]=(xyx^{-1})y^{-1}$ has $c_2$ value $-1$.
Thus the actual map from the metabelian commutator module is
\[
 c_2\bigl(A\ast[x,y]\bigr)=-A(-1,1),
 \tag{D.5.2}
\]
where $\ast$ is left conjugation. The identity follows first on
the integer linear combinations of conjugates and then on their
closure by continuity. Nakamura--Schneps use exactly
$x\mapsto x^\chi$, $y\mapsto f_\sigma^{-1}y^\chi f_\sigma$;
their equation (2.4) and Lemma 2.3 give
$f_\sigma\equiv A_\sigma\ast[x,y]$ and
$\kappa_2(\sigma)=-A_\sigma(-1,1)$
\cite[\S2, (2.1), (2.4), Lemma 2.3]{BCNakamuraSchneps2000}.
Combining these with \textup{(D.5.2)} proves
\textup{(D.5.1)}, including reduction modulo every integer.

The convention change mentioned just after their equation (2.1)
is the actual continuous free-group automorphism
\[
 \iota(x)=x^{-1},\qquad\iota(y)=y^{-1},\qquad
 \delta_2\circ\iota|_{H_2}=\iota\circ\delta_2.
 \tag{D.5.3}
\]
To verify the last equality, compare its two sides on
$x^2,y,xyx^{-1}$. The third generator is sent by $\iota$ to
$x^{-2}(xyx^{-1})^{-1}x^2$, which $\delta_2$ kills;
the first two give $x^{-1}$ and $y^{-1}$ on both sides.
Consequently $c_2(\iota f)=-c_2(f)$.
Writing $f_\sigma^{\rm Ihara}=\iota(f_\sigma)$ for the
specified change of variables gives
$-c_2(f_\sigma^{\rm Ihara})=\kappa_2(\sigma)$.
The finite distribution tests retain their literal $c_N$ maps;
all arithmetic matrix comparisons below use the displayed
Nakamura--Schneps convention.

For completeness, the automorphism $g$ associated to $(\lambda,f)$
preserves $H_N$, since it multiplies the quotient $\Z/N$ by the unit
$\lambda$. The square with $g|_{H_N}$, $g$, and the two maps
$\delta_N$ commutes. On $x^N$ both composites give $x^\lambda$;
on $y$ the extra factor $y^{c_N(f)}$ cancels in conjugation.
For $x^ayx^{-a}$, write $\lambda a=Nq+b$ with
$0<b<N$. Its image under $g$ is conjugate, by an element of $H_N$,
to $x^by^\lambda x^{-b}$; applying $\delta_N$ gives $1$.
These checks prove the square on the full profinite free basis.
\end{proof}

This opened theorem strengthens the cyclic-cover statement used in
LM7 D.20, where the displayed compatibility was used only on the
Galois image. It provides additional necessary finite tests for Task
E.3, performed below. It gives no equality between the full profinite
groups $\widehat{\mathrm{GT}}$ and $\widehat{\mathrm{GT}}_0$.

\begin{statement}{D.6}{Known}{The exact finite test and its remaining quantifier}
For a finite-index normal subgroup $N_0\triangleleft B_3$ contained in
$PB_3$, the reduction maps of finite shadow sets satisfy
\[
 \im\bigl(\widehat{\mathrm{GT}}_0\to\mathrm{GT}(N_0)\bigr)
 =\bigcap_{K\subset N_0}
 R_{K,N_0}\bigl(\mathrm{GT}(K)\bigr),
\]
where $K$ ranges over all finite-index normal subgroups of $B_3$
contained in $N_0$. This is Dolgushev--Guynee, Theorem 5.2,
Corollary 5.4 and Proposition 5.5 \cite{DE_DolgushevGuynee2024},
with the full profinite identification supplied by Harbater--Schneps,
Main Theorem (a) \cite{DE_HarbaterSchneps2000}.
\end{statement}
\begin{proof}[Verification of the inverse-limit step]
A genuine element gives a shadow in every finite quotient. Conversely,
the fibres of the reduction maps over a fixed shadow are finite sets.
A finite collection of quotient conditions is dominated by the
intersection of their kernels. Nonemptiness at every such kernel
therefore gives the finite intersection property in the product of the
finite shadow sets, and compactness produces a compatible family.
The cited theorem identifies this family, including its generation
condition, with a full profinite element. For a finite list of twelve
shadows, exclusion of each one at some quotient implies exclusion of
all twelve at the intersection of those twelve kernels. Thus the
negative conclusion admits one finite witness. Passage through any
specified collection of successful tests supplies only the conditions
for that collection.
\end{proof}

We also opened Bleher--Chinburg--Lubotzky, Lemma 2.3 and Corollary 2.4
\cite{DE_BCL2020}, and reran LM7's explicit lifting counterexample.
Those results lift automorphisms of a quotient of a finitely generated
free profinite group. LM7's finite construction uses the quotient
consisting of the free metabelian quotient together with an $A_5$
factor; the identity on the first factor and conjugation by an odd
permutation on the second lift to an automorphism. The latter sends a
five-cycle to its other $A_5$ conjugacy class. The lifted automorphism
therefore fails the required conjugacy condition on $x$ while acting
trivially on the metabelian quotient. This verifies the claimed
limitation of that particular lifting theorem. It does not rule out
another lift satisfying the peripheral equations. The missing proof
in D.1 and D.2 is exactly a solution of those equations and generation
at every finite quotient, or a finite obstruction to the proposed
extra coordinate.

\subsection{A complete finite pentagon comparison in weight nine}

\begin{statement}{D.7}{Proved}{Compatible finite three- and four-strand shadows}
Write $P_n=PB_n$, $p=101$, and
\[
 K_n=\gamma_{11}P_n\,P_n^p\triangleleft B_n\qquad(n=2,3,4),
 \qquad K_F=\gamma_{11}F\,F^p,
\]
where a power subgroup is generated by all the indicated powers.
All these groups are discrete until their finite quotients are taken.
We use $F=\langle x=A_{12},y=A_{23}\rangle\le P_3$ and the same
braid generators as in D.4. Let $\mathcal G_9$ be the finite gentle
shadows at $K_3$ with $m=0\pmod p$ and
$f\in\gamma_9(F/K_F)$. Let $\mathcal P_9\subset\mathcal G_9$ be
those satisfying the pentagon in $P_4/K_4$.
The logarithm bijection makes both sets vector spaces over
$\mathbb F_{101}$ and gives the following dimensions:
\[
\begin{array}{c|rr}
 &\mathcal G_9&\mathcal P_9\\\hline
 \log f\text{ in degrees }9,10&37&2\\
 \text{degree-nine projection}&4&1\\
 \text{degree-nine metabelian projection}&2&1.
\end{array}
\tag{D.7.1}
\]
The last right-hand space is exactly $\mathbb F_{101}G_9[x,y]$.
The induced three-strand kernel associated to $K_4$ in the finite
parenthesized-braid definition is exactly $K_3$, so
$\mathcal P_9$ consists of actual finite pentagon shadows in that
definition, not merely solutions in an unrelated Lie target.
\end{statement}
\begin{proof}
\emph{The quotient maps.}
The subgroups $K_n$ are characteristic in $P_n$, hence normal in
$B_n$. The groups $P_n/K_n$ are finitely generated nilpotent groups
of bounded finite exponent, so are finite. Every strand insertion
or cabling homomorphism $P_3\to P_4$ sends $K_3$ into $K_4$,
since homomorphisms preserve the lower central series and power
subgroups. For the inclusion of the first three strands, the
strand-forgetting retraction $P_4\to P_3$ gives
$K_4\cap P_3=K_3$. Intersecting the five preimages used in
the definition of the induced three-strand kernel therefore gives
exactly $K_3$. The analogous two-strand assertion gives
$K_2=p\mathbb Z$ and $K_{\rm ord}=p$.
The direct product $P_3=F\times\langle C_3\rangle$ gives
$K_3\cap F=K_F$.
Thus reduction induces the specified maps of finite shadow sets,
and the continuous maps from the profinite braid groups to these
finite groups are their ordinary completion maps.

\emph{Faithful logarithmic coordinates.}
The class is ten, smaller than $p$. Lazard's Theorem II.4.6,
with its homomorphism assertion, identifies the finite groups
and their Lie rings by the mutually inverse truncated logarithm
and Campbell--Hausdorff operations \cite{DE_Lazard1954}.
For the free group factor this is the free nilpotent Lie algebra
of class ten over $\mathbb F_p$. This assertion can also be read
directly from Hall collection: the logarithm of a basic
commutator is its Hall Lie word plus terms of larger degree;
this triangular matrix has diagonal one, and all denominators
through degree ten are units modulo $p$.

For four strands we keep the actual pure Artin action in these
coordinates. Put $u_i=A_{i4}$, $a=A_{12}$, $b=A_{23}$.
Forgetting strand four splits $P_4$ as
$\langle u_1,u_2,u_3\rangle\rtimes P_3$.
The center is generated by
$C_4=C_3u_1u_2u_3$; eliminating $C_3$ in the quotient by
$C_4$ gives
\[
 P_4/\langle C_4\rangle
       =F(u_1,u_2,u_3)\rtimes F(a,b).
 \tag{D.7.2}
\]
These generators and their exact conjugation relations are
\cite[(2.1)--(2.3)]{DE_CohenFalkRandell2012}.
In particular the inverse actions of $a,b$ on the free factor are
\[
\begin{array}{c|ccc}
 &u_1&u_2&u_3\\\hline
 a^{-1}(-)a&u_1u_2u_1u_2^{-1}u_1^{-1}
                  &u_1u_2u_1^{-1}&u_3\\
 b^{-1}(-)b&u_1&u_2u_3u_2u_3^{-1}u_2^{-1}
                  &u_2u_3u_2^{-1}.
\end{array}
\tag{D.7.3}
\]
The actions are trivial on the free factor's abelianization.
Their Magnus expansions raise degree by at least one; a weight-$i$
commutator of them raises degree by at least $i$. Hence
$[\gamma_iF(a,b),\gamma_jF(u_1,u_2,u_3)]
\subset\gamma_{i+j}F(u_1,u_2,u_3)$.
It follows by both inclusions that the lower central series in
(D.7.2) is the product of the two free lower central series.
Hall collection and reduction modulo $p$ therefore give faithful
coordinates consisting of the two free Lie factors. The omitted
central factor has degree one. It intersects the derived subgroup
trivially, also after this reduction, and every pentagon factor
considered below lies in degree at least nine. Thus the computation
in (D.7.2) detects the complete pentagon in $P_4/K_4$.

\emph{The complete equations are linear in the retained stratum.}
Write $X=\log x$, $Y=\log y$ and
$\log f=h_9+h_{10}$, with $h_i$ homogeneous of degree $i$.
There are $56+99=155$ Lyndon coordinates.
The two simplified finite hexagons, with $z=y^{-1}x^{-1}$, are
\[
 f(x,y)f(y,x)=1,\qquad
 f(z,x)f(y,z)f(x,y)=1.
 \tag{D.7.4}
\]
Their equivalence to the original braid hexagons for $f\in F'$
is Dolgushev--Guynee, Proposition 3.4
\cite{DE_DolgushevGuynee2024}. Here $m=0$ and
\[
 \log z=-X-Y-\tfrac12[X,Y]\pmod{\text{degree }3}.
\]
Every product of two $h_i$ has degree at least eighteen.
Consequently (D.7.4), including all terms of degree ten, is
exactly a linear system. Substitution in the two-letter tensor
algebra produces a $3072$ by $155$ matrix of rank $118$.
Its kernel has dimension $37$; its degree-nine projection has
dimension four. The program also solves for, retains, and
substitutes a degree-ten correction for each of the four leading
basis vectors.

Here are the corresponding exact four-strand substitutions.
Put $U_i=\log u_i$, $A=\log a$, $B=\log b$, and
\[
 L=U_1+U_2+U_3+\tfrac12\sum_{i<j}[U_i,U_j]
             \pmod{\text{degree }3}.
\]
If the logarithm of an inverse Artin action in (D.7.3), applied
to $U_i$, is $U_i+v_{2i}+v_{3i}$ through degree three, then
its negative logarithm is the derivation
\[
 D(U_i)=-v_{2i}-v_{3i}
             +\tfrac12D_2(v_{2i}),\qquad D_2(U_i)=v_{2i}.
 \tag{D.7.5}
\]
This follows by expanding $-\log(1+T)=-T+T^2/2$ to the
required degree. The leading parts are
\[
\begin{aligned}
 D_A(U_1)&=[U_1,U_2],&D_A(U_2)&=[U_2,U_1],&D_A(U_3)&=0,\\
 D_B(U_1)&=0,&D_B(U_2)&=[U_2,U_3],&D_B(U_3)&=[U_3,U_2].
\end{aligned}
\]
Formula (D.7.5) supplies their degree-raising-two parts directly
from the literal group words. The five pairs of arguments of the
pentagon, in right-hand-side minus left-hand-side order, are
\[
\begin{array}{c|cc}
 &\text{first logarithm}&\text{second logarithm}\\\hline
 +&A&B+U_2+\tfrac12[U_2,U_3]\\
 +&-A-L&U_3\\
 -&B&U_3\\
 -&-B-L&U_2+U_3+\tfrac12[U_2,U_3]\\
 -&A&B.
\end{array}
\tag{D.7.6}
\]
For example $A_{13}A_{23}=a^{-1}C_3$ and
$A_{12}A_{13}=C_3b^{-1}$. The degree-two corrections involving
$A$ or $B$ and $L$ vanish because both leading derivations
annihilate $U_1+U_2+U_3$. Terms of degree three in the argument
logarithms, or derivation parts raising degree by three, cannot
affect a word of weight nine through degree ten. Therefore
(D.7.5)--(D.7.6) compute the complete group equation at this
quotient, without identifying the logarithms of the pure braid
generators with homogeneous infinitesimal generators.

The pentagon's free two-generator component cancels. Its free
three-generator component has $3^9$ and $3^{10}$ tensor
coefficients. The degree-nine equations reduce the $37$-dimensional
hexagon space to dimension $34$; all $59049$ degree-ten equations
then reduce it to dimension two. Projecting this kernel to degree
nine gives rank one, and applying the metabelian quotient gives
the asserted gamma line. All surviving basis vectors were
substituted into every hexagon and pentagon coefficient.

\emph{Generation.}
The map on $P_3/K_3$ induced by a retained pair sends
$x$ to $x$, $y$ to $f^{-1}yf$, and $C_3$ to $C_3$.
It is the identity on the Frattini quotient of this finite
$p$-group, and is therefore surjective. The two-strand map is
also surjective. Definition 2.9 and Proposition 2.10 of
Dolgushev--Le--Lorenz \cite{DE_DolgushevLeLorenz2024} now
give the generation condition on four strands for every retained
pentagon pair. The same three-strand argument gives the gentle
generation condition. This proves that the computed solutions
are the claimed finite shadows.

The executable sources are
\texttt{de\_gentle\_weight9.py} and
\texttt{de\_pentagon10\_finite.py}; their JSON outputs retain
the basis logarithms, and their compressed certificates retain
the full coefficient matrices. The latter run completed in
$501.307$ seconds. Its assertions verify every surviving
coefficient, the projection ranks, and equality with the gamma
line. No search over a subset of the $155$-coordinate stratum
is used.
\end{proof}

\begin{statement}{D.8}{Proved}{An integer pentagon certificate and the weight-nine comparison}
Let $k$ be a field of characteristic zero or characteristic at
least five. In the free Lie algebra $L_k(X,Y)$ put
\[
 V_9=\{h\in L_k(X,Y)_9:
 h(X,Y)+h(Y,X)=0,\quad
 [X,h(X,Y)]=[X+Y,h(-X-Y,Y)]\}.
\]
Then $\dim V_9=4$, its metabelian image has dimension two,
and the infinitesimal pentagon map $P:V_9\to(\mathfrak t_4)_9$
has rank three. Its one-dimensional kernel has metabelian image
$k\,G_9(X,Y)[X,Y]$. The complete finite projection in D.7 is this kernel
when $k=\mathbb F_{101}$.
For every $\ell\ge5$, applying the corresponding rational
leading-term map to the full profinite pentagon equation and
using D.20 gives precisely $\mathbb Q_\ell G_9(X,Y)[X,Y]$ as the
weight-nine metabelian image of $\widehat{\rm GT}$ and of
the Galois image.
\end{statement}
\begin{proof}
The following exact integer certificate also explains the
comparison with the finite calculation. Generate the $56$
Lyndon words of length nine lexicographically, and use standard
bracketing: split at the longest proper Lyndon suffix, then
bracket the two recursively bracketed factors. Expand in the
free associative algebra. The two equations defining $V_9$
give an integer $1536$ by $56$ matrix $C$. Its $52$ nonzero
Smith factors are all one. A four-column integer kernel basis
$H$ retained in the JSON certificate has Smith factors
$1,1,2,24$. Hence it is a basis of the kernel over every field
in the statement, not merely over $\mathbb Q$.

We evaluate the pentagon in the faithful Lie model
\[
 \mathfrak t_4/(\text{degree-one center})
       =L(U_1,U_2,U_3)\rtimes L(A,B),
\]
with the leading derivations displayed in D.7. It is the
associated graded of (D.7.2); the integral braid relations are
also recorded in \cite[(4.1)]{DE_CohenFalkRandell2012}.
For the five arguments take the degree-one parts of (D.7.6).
All eighteen infinitesimal braid relations are checked by the
program, as is the cancellation of the free $L(A,B)$ component.
The three tensor rows of $P(H)$ labelled
$111111223$, $111111232$, $111121123$ are
\[
 \begin{pmatrix}
 -6&-6&6&-5\\
 254&510&-446&-1291\\
 -150&-302&246&700
 \end{pmatrix}.
 \tag{D.8.1}
\]
The determinant of the first three columns is
$27648=2^{10}3^3$. The vector
\[
                 v=(-3299,1326,-1613,432)^{\mathsf T}
 \tag{D.8.2}
\]
annihilates all $19683$ tensor coefficients of $P(H)$, verified
with integer arithmetic. Its coordinates have greatest common
divisor one. Thus $P$ has rank exactly three and kernel $kv$
in each of the stated characteristics.

In the ordered monomial coordinates
$X^iY^{7-i}[X,Y]$ for $0\le i\le7$, the metabelian image of
$H$ is
\[
 \begin{pmatrix}
 0&0&0&-2\\
 0&0&0&-8\\
 2&2&-2&-17\\
 6&6&-6&-23\\
 6&6&-6&-23\\
 2&2&-2&-17\\
 0&0&0&-8\\
 0&0&0&-2
 \end{pmatrix}.
 \tag{D.8.3}
\]
Its nonzero Smith factors are $1,4$. Applied to (D.8.2), it
gives
\[
 -96(9,36,84,126,126,84,36,9)^{\mathsf T}=-96G_9.
\]
This proves all the dimension and coordinate assertions.
The programs \texttt{de\_pentagon9.py} and
\texttt{de\_pentagon9\_rational.py} construct the matrices
from the Lie relations; the second performs the integer and
rational calculations anew, including both Smith forms and
every zero coefficient.

For the final comparison let a full profinite element have
weight-nine leading action in the filtration of Part D.
Its $\ell$-adic multiplier is one. In the normalized gauge
the action fixes $x$ and its conjugator belongs to $F'$.
A lower-degree inner action fixing $x$ would have an inner
generator centralizing $X$. In degree at least two this is
zero in the free Lie algebra; the degree-one possibility is
excluded by the derived-subgroup normalization. The conjugator
therefore has no term below degree nine. Its leading word
belongs to $V_9$ by the peripheral and symmetry equations.

The insertion homomorphisms of pure braid groups preserve
their lower central series. Their associated graded maps
send each pure generator to the corresponding sum of
infinitesimal generators; this follows by abelianizing the
explicit product formulas in the five insertion maps.
In degree nine the product of pentagon factors becomes their
sum, since products of two degree-nine terms have degree
eighteen. The leading word consequently lies in the exact
kernel just computed over $\mathbb Q_\ell$. This is the
upper inclusion, with its source and target specified.
The actual Galois construction in D.20 supplies a nonzero
element in the same filtered full profinite source, giving
the reverse inclusion on rational spans. This derives R9 in
weight nine from the certified pentagon map.

Finally, other multipliers act nontrivially in degree one,
and lower nonzero leading words belong to lower filtration
pieces. They do not define a weight-nine leading class under
the map in Part D. This is why D.7 enumerates precisely the
stratum needed for this comparison. The quotient detects
an extra finite gentle direction, but the calculation does
not construct a compatible family in the full profinite
gentle group. The remaining existence question is D.2.
\end{proof}

